\section{Experiments}
\label{sec:experiments}

We evaluate OPTS-TTPO on verifiable mathematical reasoning, discrete control, and continuous control. Unless otherwise stated, all comparisons use a matched-budget protocol: LLM experiments are matched by the number of completed answers, while Atari and MuJoCo experiments are matched by environment interaction steps. Additional implementation details are summarized in Appendix~\ref{app:experiment-setup}, Table~\ref{tab:experiment-setup}.

\subsection{LLM Training-Time Search}
\label{sec:llm-train-search}

\subsubsection{Setup}

We train \texttt{Qwen3-1.7B} with the VeRL framework on verifiable mathematical reasoning. The training data consists of \texttt{math12k} and the competition subset of \texttt{NuminaMath-1.5-RL-Verifiable}. We evaluate on \texttt{math12k} test (\texttt{MATH500}), \texttt{minervamath}, \texttt{amc23}, and \texttt{aime25}. Baselines include PPO, DAPO, GPG, REINFORCE++, and OPTS-TTPO.

This experiment asks whether training-time search improves the learned policy under the same answer budget, and whether the gains appear consistently in training dynamics rather than only at a single checkpoint.

\subsubsection{Pending Results and Target}

This experiment is currently running. The final version will report validation curves (\texttt{avg@n}, \texttt{pass@n}, and \texttt{cons@n}), a final-checkpoint table with \texttt{avg@32}/\texttt{pass@32}/\texttt{cons@32}, and the checkpoint selection rule. The goal is to determine whether OPTS-TTPO improves average sample quality or correct-solution coverage over PPO and DAPO.

\subsection{LLM Test-Time Search}
\label{sec:llm-test-search}

We also instantiate OPTS as a test-time search method. The inference budget is matched to repeated-sampling baselines, but OPTS adaptively allocates continuation budget according to prefix quality rather than sampling complete answers independently.

\subsubsection{Pending Results and Target}

This experiment is currently running. The final version will report \texttt{avg@k}/\texttt{pass@k}/\texttt{cons@k} curves for $k\in\{8,16,32,64,128\}$, and compare against i.i.d. sampling and self-consistency under the same completed-answer budget.

\subsection{Cross-Domain Validation: Atari and MuJoCo}
\label{sec:cross-domain}

To test whether OPTS-TTPO transfers beyond language reasoning, we compare it with PPO on Atari-57 discrete control and five MuJoCo continuous-control tasks. These domains differ from LLM reasoning in observation modality, action space, and budget unit, making them a useful testbed for whether OPTS-TTPO acts as a general rollout-budget reallocation mechanism.

\subsubsection{Atari-57}
\label{sec:atari}

Atari-57 has highly heterogeneous reward scales across games. We therefore report task-level win counts in the main text and provide the full learning curves in Appendix~\ref{app:atari-curves}, Figure~\ref{fig:atari-all-curves}. For each game, we average aligned logs across three random seeds and compare two summary statistics: the mean return over the full training run and the mean return over the last 100 logged points. Table~\ref{tab:atari-win-counts} reports the number of games won by PPO and OPTS-TTPO.

\begin{table}[t]
\centering
\small
\begin{tabular}{lcc}
\hline
Metric & PPO & OPTS-TTPO \\
\hline
Full-training mean return & 26 & \textbf{31} \\
Last-100-log mean return & 24 & \textbf{33} \\
\hline
\end{tabular}
\caption{Task-level win counts on Atari-57. Each game contributes one win to the method with the higher score under the corresponding metric.}
\label{tab:atari-win-counts}
\end{table}

OPTS-TTPO obtains more wins under both metrics, suggesting that on-policy rebranching can scale to high-dimensional visual control. The gain in full-training mean return indicates that rebranching can improve the learning process across many tasks, while the gain in the tail metric suggests that the effect is not limited to early exploration. At the same time, PPO still wins on a substantial subset of games, showing that the benefit is task-dependent rather than a monotonic improvement on every reward structure.

\subsubsection{MuJoCo}
\label{sec:mujoco}

We evaluate on \texttt{Hopper-v4}, \texttt{Walker2d-v4}, \texttt{HalfCheetah-v4}, \texttt{Ant-v4}, and \texttt{Humanoid-v4}. Figure~\ref{fig:mujoco-curves} shows the latest learning curves, with OPTS-TTPO in red and PPO in blue.

\begin{figure*}[t]
\centering
\includegraphics[width=\textwidth]{figures/all_tasks_mujoco.png}
\caption{Learning curves on MuJoCo. Red denotes OPTS-TTPO and blue denotes PPO; shaded regions show variation across random seeds.}
\label{fig:mujoco-curves}
\end{figure*}

The curves show clear gains for OPTS-TTPO on \texttt{HalfCheetah-v4} and \texttt{Humanoid-v4}, and a stronger late-stage trend on \texttt{Walker2d-v4}. On \texttt{Hopper-v4}, the two methods are close, with PPO slightly ahead near the end. \texttt{Ant-v4} exhibits large seed variance, making a stable advantage difficult to establish. Consistent with the Atari results, MuJoCo suggests that rebranching can improve learning speed or final return in some continuous-control tasks, but its benefit depends on the task dynamics and the selected branching positions.

\subsection{Policy-Gradient Variance Reduction}
\label{sec:variance-reduction}

Finally, we test the variance-reduction mechanism behind OPTS-TTPO. For each MuJoCo task, we use a fixed checkpoint trained by PPO for 1M steps and approximate a reference actor gradient $g^\star$ with a large PPO rollout buffer. For a mini-batch of size $B$, we estimate $\|\hat g_B-g^\star\|_2^2$ by bootstrap sampling and use it as a policy-gradient variance proxy; lower values indicate more stable gradient estimates.

\subsubsection{High-Advantage Samples Have Lower Variance}

Figure~\ref{fig:pg-pos-neg-variance} compares samples selected by step-level GAE advantage. Across all five MuJoCo tasks, high-advantage samples have lower policy-gradient variance for almost all batch sizes. This supports the design of OTRC, which allocates budget toward high-advantage regions. Variance decreases with batch size for both groups, but the persistent gap shows that sample selection still affects gradient quality.

\begin{figure*}[t]
\centering
\includegraphics[width=\textwidth]{figures/pg_variance_verify1.png}
\caption{Policy-gradient variance for high- and low-advantage samples. Red denotes high-advantage samples and blue denotes low-advantage samples; the vertical axis is on a log scale, where lower is better.}
\label{fig:pg-pos-neg-variance}
\end{figure*}

\subsubsection{OPTS-TTPO Reduces Matched-Batch Gradient Variance}

Figure~\ref{fig:pg-scaling-variance} compares PPO rollout samples and OPTS-TTPO tree samples under the same policy checkpoint and matched batch sizes. OPTS-TTPO uses the same $1/W$ inverse-probability weighting as in training to correct for the sampling-measure change induced by rebranching. OPTS-TTPO is generally below PPO on \texttt{Hopper-v4}, \texttt{Walker2d-v4}, \texttt{HalfCheetah-v4}, and \texttt{Ant-v4}, indicating lower-variance gradient estimates at the same batch size. \texttt{Humanoid-v4} is the main exception, suggesting that variance reduction also depends on task dynamics and branching-position selection.

\begin{figure*}[t]
\centering
\includegraphics[width=\textwidth]{figures/scaling_variance_verify2.png}
\caption{Policy-gradient variance scaling for PPO and OPTS-TTPO under matched batch sizes. Red denotes OPTS-TTPO and blue denotes PPO; the vertical axis is on a log scale, where lower is better.}
\label{fig:pg-scaling-variance}
\end{figure*}
